\section{Relation to numerical automata and positive realization}\label{sec:comparison54}
The input/output object and the order of quantifiers are decisive in comparing realization theorems. Classical probabilistic automata assign numerical acceptance probabilities to words, while stochastic-language results also study threshold recognition \cite{Rabin,Paz}. A numerical equivalence of two fixed machines is not the assertion that for every horizon there exists a new bounded-width approximate machine. Conversely, language equivalence need not preserve the numerical probabilities.

For weighted automata over a field, the finite-Hankel-rank characterization gives finite linear representations and identifies their minimum dimension \cite{BR}. It does not enforce nonnegative rows, stochastic normalization, or simplex-valued decoders. Our common-row contraction concerns these constraints simultaneously. The invariant observable quotient in Theorem~\ref{thm:observable53} uses the same reachable/invisible linear reduction; its additional conclusion is a finite-group criterion for exact positive realization with the full horizon-dependent quantifier.

The classical positive-realization problem asks for a positive system representing specified input/output data, and distinguishes positive dimension from unrestricted linear dimension \cite{BF04}. Invariant cones and polyhedral sections are part of that theory, not new notions introduced here. Proposition~\ref{prop:polygon54} is an explicit common invariant-enclosure construction. The new matched statement pairs it with a full-subcritical lower estimate against all clocked stochastic machines. Neither one cutwise nonnegative factorization nor a finite-rank Hankel matrix supplies that lower estimate.

Brodsky and Pippenger identify bounded-error measure-once quantum languages as group languages \cite[Theorem 3.3]{BP}. Their probabilistic simulation result preserves a cutpoint language with a potentially changed cutpoint \cite[Theorem 3.6]{BP}; it is not uniform additive preservation of every acceptance probability. Our categorical unitary corollary fixes an entire terminal measurement law in total variation and allows horizon-dependent stochastic approximants. The distinction is already visible in the $d$-outcome threshold $1-1/d$.

The closure theorem and strict-threshold decision method of Blondel, Jeandel, Koiran, and Portier \cite{BJKP}, and the effective Zariski-closure algorithms \cite{DJK,NPSW}, are algorithmic premises of Section~\ref{sec:effective54}. Computing the closure is not a new claim of this article. The additional outputs are a constrained minimax noise threshold, an attaining categorical decoder, and a terminating construction of an all-machine occupation constant. Recent Lie-algebra-based closure computations \cite{deGraaf} give an alternative representation of the algebraic closure, but their complex connected components still require real-component analysis for the present application.

Peter--Weyl approximation, Haar averaging, and the representative-function algebra are classical compact-group and almost-periodic-function machinery \cite{Folland}. Uniform approximation by a finite sum of matrix coefficients gives a finite \emph{linear} representation of an output function. It does not give positive stochastic rows or the repeated information loss quantified by \eqref{eq:occupation54}. Similarly, Steinberg's minimal compact topological automaton and enveloping-monoid results \cite[Theorems 2.2 and 2.6]{Steinberg} concern recognition of a language. His compact topological-monoid criterion \cite[Proposition 2.8]{Steinberg} uses clopen recognition, not a continuous simplex-valued response with a prescribed error norm. The finite-quotient radius formula is specific to groups. The stationarization theorem itself now holds for irreversible topological monoids with eventual approximate returns (Theorem~\ref{thm:stationarization56}).

Distribution-based bisimulation metrics provide robust comparisons for specified probabilistic automata \cite{FSZ}. They do not by themselves impose a minimum over all horizon-specific stochastic realizations of an external group experiment. The finite future-output equivalence in Theorem~\ref{thm:quotient54} is deliberately a component-quotient statement, rather than a claim that approximate equivalence always has a unique minimal positive stochastic realization.

\begin{center}\small
\begin{tabular}{p{.27\textwidth}p{.31\textwidth}p{.32\textwidth}}
\toprule
Comparison object & Error or equivalence & Machine quantifier\\
\midrule
Compact stochastic groups \cite{Flor} & Idempotent and permutation structure; no external error specification & One matrix group or semigroup kernel\\
Advice regular languages \cite{FP} & Neutral-letter recognition & One-scan advice automata; Boolean output\\
Weighted series \cite{BR} & Exact scalar coefficients; linear rank & One linear representation\\
Group-language simulation \cite{BP} & Cutpoint language; not additive probability error & One automaton\\
Positive realization \cite{BF04} & Exact specified numerical response & One positive system\\
Topological recognition \cite{Steinberg} & Clopen language acceptance & One recognizing action\\
Bisimulation metrics \cite{FSZ} & Distributional comparison of given systems & Two specified automata\\
Theorem~\ref{thm:stationarization56} & Same closed numerical error and label bound & Arbitrarily many narrow cuts; one extracted stationary machine\\
Theorem~\ref{thm:radius54} & Worst-word norm error; TV for a whole categorical law & A new clocked stochastic machine at every horizon\\
Theorem~\ref{thm:matched54} & Binary TV, every fixed $\varepsilon<\rho/2$ & Same nonuniform lower model; exact upper construction\\
\bottomrule
\end{tabular}
\end{center}
The contribution combines stationary purification and its exact finite-quotient criterion with a sharp constrained clocked radius, phasewise arbitrary-width occupation, and, for the specified Diophantine family, a matching full-subcritical growth law. The elementary radius center, the arithmetic pigeonhole argument, the closure algorithm, and the harmonic-analysis ingredients are not asserted individually new. The accompanying theorem-level audit records the closest comparisons and their changed hypotheses; it is not a claim of exhaustive priority clearance.


\subsection{The classical matrix-semigroup boundary}
Flor's description of nonnegative matrix groups \cite[Theorems 2 and 3]{Flor} identifies their idempotent geometry and the finiteness of bounded groups; his account explicitly credits Brown and Schwarz for earlier compact and stochastic cases. We do not claim this matrix structure as new. Theorem~\ref{thm:purification55} uses a minimal-rank idempotent in a joint semigroup, not just in the matrix projection: its physical coordinate must be the group identity. Sandwiching each command then preserves all uniform error inequalities by closure, even when the original stochastic rows violate every nontrivial physical group relation. The resulting initialization may be random, a point essential to exact state minimization.

The stationary theorem has no return hypothesis. The clocked theorem has a different quantifier and needs synchronization; the empty-return example and the phase-dependent example show why the distinction cannot be removed by terminology. The exact-return-length reduction is an additive-semigroup argument, whereas the profinite zero-error converse uses finite Hankel rank and continuous finite-dimensional representations. The new stationarization theorem removes the reversibility requirement for the equality of limiting width and stationary width, but does not supply a finite-quotient radius formula for an arbitrary irreversible monoid.

The finite-quotient existence theorem does not say that $[H:U]$ is the optimal number of stochastic labels. The actual minimum is the permutation-feasibility invariant in Theorem~\ref{thm:stationaryminimum55}; the finite cyclic example separates them. For infinitely many components, the infimum of finite-quotient radii can fail to be attained, even at zero. Finally, the exactly matched quantitative law retains its specified Diophantine alphabet. Theorem~\ref{thm:spherical56} gives the retained block estimate, strengthened to a matching power law by Theorem~\ref{thm:sharpoccupation61}, for spherical representations with a spectral gap; Corollary~\ref{cor:SO356} gives a concrete rational noncommuting instance. The exact unit-amplitude endpoint is not covered by these upper constructions. These distinctions delimit the proved statements without weakening any of the retained conclusions.

\subsection{What the reductions add}
The algebraic ingredients of Lemmas~\ref{lem:idempotent55} and \ref{lem:simplex55} are the classical minimum-rank kernel mechanism, the group corner of a compact matrix semigroup, and the recurrent simplex of a stochastic idempotent. Flor's Theorem~2 describes nonnegative idempotents; his Theorem~3 and preceding Lemma~3 give the permutation structure of bounded nonnegative groups \cite{Flor}. The paper explicitly attributes earlier cases to Brown and Schwarz. The present proof adds the physical coordinate to the closure, chooses the idempotent over its identity, and tracks the closed error balls through sandwiching. Theorem~\ref{thm:equivariant56} further averages the finite kernel above the physical identity and descends to an actual continuous $H$-action. These are specified realization reductions using classical structure, not a new classification of matrix-semigroup kernels.

For deterministic Boolean recognition, neutral-letter removal of advice is an established phenomenon. Fijalkow--Paperman \cite[Theorem~10]{FP} prove the Crane Beach property for advice regular languages by padding distinguishing words and suffixes, and their account relates one-scan programs to earlier Ramsey arguments. Their Myhill--Nerode argument in fact preserves the deterministic distinguishability bound $K$; preservation of a state bound alone is not claimed as new here. A stochastic $k$-label register has a continuum of state distributions, so this discrete class count does not apply. Theorem~\ref{thm:stationarization56} instead colors approximate tuples of stochastic composite kernels, extracts one common tuple, and absorbs a final filler matrix into the decoder. It preserves the same numerical tolerance at equality and the same $k$, and bounds all narrow-cut occurrences with unrestricted intermediate widths. Eventual approximate returns further replace a literal neutral letter. The finite Ramsey theorem itself is classical \cite{Ramsey}.

The spherical lower estimate uses a classical spectral gap, cap measure, and smoothing; its upper estimate uses a polytope enclosing a contracted sphere. The contribution is their realization-theoretic combination with the all-profile conditional-centroid bound. The rational $\mathrm{SO}(3)$ example invokes \cite{BG2012} as a deep external input. Similarly, the lower-complexity input in Theorem~\ref{thm:ETR56} is Shitov's nonnegative-factorization universality theorem \cite{Shitov}; the stochastic normalization and finite-group existential upper reduction are given explicitly. Neither external theorem is claimed as part of our proof technology's novelty.

\subsection{Companion results}
The companion supplement gives additional results on compatible reachable sections, one-surplus geometry, stable sections, finite-horizon certificates, word distortion, two-state duality, finite-bit compilation, and global residual hierarchies. The scalar two-extreme argument and earlier quantitative constants are also retained there. The present main proofs are self-contained apart from the classical results explicitly cited; none of these companion derivations is used to conceal a missing step in a principal theorem.

\subsection{Orbit geometry, positive matrix outputs, and recognition}
The homogeneous-space covering estimates and the smooth orbit geometry used in Section~\ref{sec:uniform57} are classical; see, for example, Szarek~\cite{Szarek57}. The spectral-gap input is Bourgain--Gamburd~\cite{BG2012}. Dense algebraic free generators in Section~\ref{sec:extreme57} come from Breuillard--Gelander~\cite[Theorem~1.1]{BreuillardGelander57}; the later repair and supplementary proof details are recorded in~\cite{BreuillardGelanderCorrection57}. None of these results is claimed as a new group-theoretic theorem here. The realization argument requires, in addition, a contraction uniform over every conditional-centroid direction, one common positive row for each command and label, and a terminal decoder in the original legal matrix set.

Quantum finite automata already exhibit strong state-complexity separations in language recognition; Ambainis--Freivalds~\cite{AF57} is a relevant early comparison. Theorem~\ref{thm:freeendpoint57} instead asks a classical stochastic machine to reproduce one entire legal density-matrix vector at every input word, permits horizon-specific redesign and a free clock, and compares exact and fixed-error width for the same interface. Its exact lower bound uses extreme points of the decoder set; it does not apply unchanged to a single acceptance probability or to an enlarged unconstrained output simplex. A finite tomography map preserves width only when its legal image and reconstructed norm are retained. This comparison identifies different quantifiers and resources, not an exhaustive priority claim about all quantum or probabilistic automata formalisms.

The spectral-entropy argument removes the logarithm from the nonabelian lower bound. Equality of lower and upper powers requires equality of least-orbit and target-orbit dimensions; no universally matching law for every representation is asserted. General online outputs, adaptive feedback, and unrestricted compact semigroup radius formulas require additional arguments. The same-width stationarization result for topological monoids remains available under its actual approximate-return hypothesis; it does not provide that missing radius formula.
